\documentclass[a4paper,12pt]{article}
\usepackage[utf8]{inputenc}
\usepackage[T1]{fontenc}
\usepackage[french]{babel}
\usepackage{geometry}
\usepackage{graphicx}
\usepackage{tikz}
\usepackage{xcolor}
\usepackage{fontawesome5}
\usepackage{tcolorbox}
\usepackage{booktabs}
\usepackage{array}
\usepackage{multirow}

\geometry{margin=2.5cm}
\usetikzlibrary{shapes.geometric, arrows, positioning, shadows}

% Couleurs personnalisées
\definecolor{primaryblue}{RGB}{34, 197, 94}
\definecolor{accentgreen}{RGB}{16, 185, 129}
\definecolor{alertred}{RGB}{239, 68, 68}
\definecolor{darkbg}{RGB}{15, 23, 42}
\definecolor{lightgray}{RGB}{243, 244, 246}

\title{\textbf{Smart Shepherd - Vue d'Ensemble du Dashboard}}
\author{Interface d'Administration IoT}
\date{\today}

\begin{document}

\maketitle

\begin{tcolorbox}[colback=primaryblue!10, colframe=primaryblue, title={\faIcon{tachometer-alt} \textbf{Tableau de Bord Intelligent}}]
Le dashboard Smart Shepherd fournit une surveillance en temps réel complet du troupeau avec des métriques de performance, des alertes intelligentes et des visualisations de données interactives.
\end{tcolorbox}

\section{Architecture du Dashboard}

\begin{center}
\begin{tikzpicture}[
    node distance=2cm,
    box/.style={rectangle, rounded corners, draw=primaryblue, fill=lightgray, text width=3cm, minimum height=1cm, text centered, font=\small},
    metric/.style={rectangle, rounded corners, draw=accentgreen, fill=accentgreen!20, text width=2.5cm, minimum height=0.8cm, text centered, font=\footnotesize},
    alert/.style={rectangle, rounded corners, draw=alertred, fill=alertred!20, text width=2.5cm, minimum height=0.8cm, text centered, font=\footnotesize}
]

% Header
\node[box] (header) {En-tête de Navigation\\Contrôles Utilisateur};
\node[box, below=of header] (main) {Zone Principale\\Grid Layout};

% KPI Cards
\node[metric, below left=0.5cm and 1cm of main] (total) {Total Troupeau\\{\Large\faIcon{users}}};
\node[metric, below=of total] (active) {En Ligne\\{\Large\faIcon{wifi}}};
\node[metric, below=of active] (alerts) {Alertes\\{\Large\faIcon{exclamation-triangle}}};
\node[metric, below=of alerts] (outofzone) {Hors Zone\\{\Large\faIcon{map-marker-alt}}};

% Charts
\node[box, below right=0.5cm and 1cm of main] (charts) {Graphiques Temps Réel\\Métriques d'Activité};
\node[alert, below=of charts] (alertfeed) {Flux d'Alertes\\Notifications Critiques};

% Weather Widget
\node[box, above right=0.5cm and 1cm of main] (weather) {Widget Météo\\Conditions Actuelles};

% Connections
\draw[->, thick, primaryblue] (header) -- (main);
\draw[->, thick, accentgreen] (main) -- (total);
\draw[->, thick, accentgreen] (main) -- (active);
\draw[->, thick, accentgreen] (main) -- (alerts);
\draw[->, thick, accentgreen] (main) -- (outofzone);
\draw[->, thick, primaryblue] (main) -- (charts);
\draw[->, thick, alertred] (charts) -- (alertfeed);
\draw[->, thick, primaryblue] (main) -- (weather);

\end{tikzpicture}
\end{center}

\section{Composants Principaux du Dashboard}

\subsection{Cartes KPI (Key Performance Indicators)}

\begin{table}[h]
\centering
\begin{tabular}{|l|c|c|l|}
\hline
\textbf{Indicateur} & \textbf{Valeur} & \textbf{Icône} & \textbf{Description} \\
\hline
Total Troupeau & \textcolor{primaryblue}{\textbf{Variable}} & \faIcon{users} & Nombre total de moutons équipés \\
En Ligne & \textcolor{accentgreen}{\textbf{Variable}} & \faIcon{wifi} & Appareils connectés \\
Alertes & \textcolor{alertred}{\textbf{Variable}} & \faIcon{exclamation-triangle} & Notifications actives \\
Hors Zone & \textcolor{alertred}{\textbf{Variable}} & \faIcon{map-marker-alt} & Animaux hors périmètre \\
\hline
\end{tabular}
\caption{Indicateurs de performance en temps réel}
\end{table}

\subsection{Métriques Secondaires}

\begin{tcolorbox}[colback=lightgray, colframe=primaryblue, title={\faIcon{chart-line} Métriques de Santé}]
\begin{itemize}
    \item \textbf{Niveau Batterie :} Pourcentage moyen avec alertes < 20\%
    \item \textbf{Température Troupeau :} Moyenne 38.5°C (normale)
    \item \textbf{Charge Ingestion :} Flux de télémétrie en pourcentage
\end{itemize}
\end{tcolorbox}

\subsection{Graphiques d'Activité}

\begin{center}
\begin{tikzpicture}[scale=0.8]
    % Axes
    \draw[->] (0,0) -- (8,0) node[anchor=north] {Temps};
    \draw[->] (0,0) -- (0,4) node[anchor=east] {Valeur};
    
    % Grid
    \foreach \x in {1,2,3,4,5,6,7}
        \draw[gray!30, dashed] (\x,0) -- (\x,4);
    \foreach \y in {1,2,3}
        \draw[gray!30, dashed] (0,\y) -- (8,\y);
    
    % Data lines
    \draw[primaryblue, ultra thick] plot coordinates {(0.5,1) (1.5,1.5) (2.5,2) (3.5,2.3) (4.5,2.8) (5.5,3.2) (6.5,3.5) (7.5,3.8)};
    \draw[alertred, ultra thick, dashed] plot coordinates {(0.5,0.5) (1.5,0.8) (2.5,1.2) (3.5,1.0) (4.5,1.5) (5.5,1.8) (6.5,2.0) (7.5,2.2)};
    
    % Legend
    \node[primaryblue] at (6,3.5) {Troupeau Actif};
    \node[alertred] at (6,3.0) {Niveau Alertes};
    
    % Labels
    \foreach \x/\label in {1/T1, 2/T2, 3/T3, 4/T4, 5/T5, 6/T6, 7/T7}
        \node[below] at (\x,0) {\tiny\label};
\end{tikzpicture}
\end{center}

\section{Système d'Alertes}

\subsection{Types d'Alertes}

\begin{table}[h]
\centering
\begin{tabular}{|l|c|c|l|}
\hline
\textbf{Type} & \textbf{Sévérité} & \textbf{Couleur} & \textbf{Action} \\
\hline
Température Élevée & CRITIQUE & \textcolor{alertred}{Rouge} & Intervention immédiate \\
Batterie Faible & WARNING & \textcolor{orange}{Orange} & Recharge planifiée \\
Hors Zone & WARNING & \textcolor{orange}{Orange} & Vérification géographique \\
Fréquence Cardiaque & CRITIQUE & \textcolor{alertred}{Rouge} & Examen vétérinaire \\
Inactivité & INFO & \textcolor{blue}{Bleu} & Monitoring continu \\
\hline
\end{tabular}
\caption{Classification des alertes par sévérité}
\end{table}

\subsection{Flux d'Alertes Temps Réel}

\begin{tcolorbox}[colback=alertred!10, colframe=alertred, title={\faIcon{bell} Notifications Critiques}]
\begin{itemize}
    \item \textbf{Alertes Non Lues :} Compteur en temps réel
    \item \textbf{Filtrage par Sévérité :} CRITIQUE > WARNING > INFO
    \item \textbf{Historique :} 100 dernières alertes conservées
    \item \textbf{Actions Rapides :} Marquer comme lue, Détails, Supprimer
\end{itemize}
\end{tcolorbox}

\section{Contrôles et Interactions}

\subsection{Boutons de Contrôle Principaux}

\begin{center}
\begin{tikzpicture}[
    button/.style={rectangle, rounded corners, draw=primaryblue, fill=primaryblue!20, text width=3cm, minimum height=0.8cm, text centered}
]

% Simulation Toggle
\node[button] (sim) at (0,0) {\faIcon{play} Simulation};
\node[below=0.1cm of sim, font=\tiny] {Basculer mode};

% Pause/Resume
\node[button] at (4,0) {\faIcon{pause} Pause};
\node[below=0.1cm of 4,0, font=\tiny] {Mettre en pause};

% Connection Status
\node[button] at (8,0) {\faIcon{wifi} Connecté};
\node[below=0.1cm of 8,0, font=\tiny] {Statut MQTT};

\end{tikzpicture}
\end{center}

\subsection{État de Connexion}

\begin{table}[h]
\centering
\begin{tabular}{|l|c|c|}
\hline
\textbf{État} & \textbf{Indicateur} & \textbf{Description} \\
\hline
MQTT Connecté & \textcolor{accentgreen}{\faIcon{check-circle}} & Communication en temps réel active \\
Mode Hors-ligne & \textcolor{orange}{\faIcon{exclamation-triangle}} & Données locales utilisées \\
Passerelle Hors-ligne & \textcolor{alertred}{\faIcon{times-circle}} & Aucune connexion disponible \\
\hline
\end{tabular}
\caption{États de connexion et leurs indicateurs visuels}
\end{table}

\section{Widget Météo}

\begin{tcolorbox}[colback=lightgray, colframe=primaryblue, title={\faIcon{cloud-sun} Conditions Météorologiques}]
\begin{itemize}
    \item \textbf{Température Actuelle :} Affichage en temps réel
    \item \textbf{Conditions :} Ensoleillé, Nuageux, Pluvieux
    \item \textbf{Prévisions :} 24 heures suivantes
    \item \textbf{Impact Troupeau :} Recommandations basées sur la météo
\end{itemize}
\end{tcolorbox}

\section{Performance et Optimisation}

\subsection{Métriques de Performance}

\begin{table}[h]
\centering
\begin{tabular}{|l|c|c|l|}
\hline
\textbf{Métrique} & \textbf{Valeur Cible} & \textbf{Actuel} & \textbf{Statut} \\
\hline
Temps de Chargement & < 2s & Variable & \textcolor{orange}{À optimiser} \\
Mises à Jour Chart & 5s & 5s & \textcolor{accentgreen}{OK} \\
Mémoire Utilisée & < 100MB & Variable & \textcolor{orange}{À surveiller} \\
Réactivité UI & < 100ms & Variable & \textcolor{orange}{À améliorer} \\
\hline
\end{tabular}
\caption{Indicateurs de performance du dashboard}
\end{table}

\subsection{Optimisations Implémentées}

\begin{itemize}
    \item \textbf{Lazy Loading :} Chargement différé des composants lourds
    \item \textbf{Memoization :} Mémorisation des calculs coûteux
    \item \textbf{Virtual Scrolling :} Pour listes longues d'alertes
    \item \textbf{Code Splitting :} Division du bundle JavaScript
    \item \textbf{Caching :} Mise en cache des données fréquemment accédées
\end{itemize}

\section{Accessibilité et UX}

\subsection{Caractéristiques d'Accessibilité}

\begin{itemize}
    \item \textbf{Navigation Clavier :} Accès complet via tabulation
    \item \textbf{Contrastes Élevés :} Respect WCAG 2.1 AA
    \item \textbf{Lecteurs Écran :} Labels ARIA complets
    \item \textbf{Mode Sombre :} Réduction de la fatigue oculaire
    \item \textbf{Responsive Design :} Adaptation mobile/tablette
\end{itemize}

\subsection{Expérience Utilisateur}

\begin{tcolorbox}[colback=primaryblue!10, colframe=primaryblue, title={\faIcon{user-circle} UX Design}]
\begin{itemize}
    \item Interface intuitive avec feedback visuel immédiat
    \item Animations fluides et transitions naturelles
    \item Informations hiérarchisées par importance
    \item Actions rapides accessibles en un clic
    \item État système toujours visible
\end{itemize}
\end{tcolorbox}

\section{Conclusion}

Le dashboard Smart Shepherd représente une solution de monitoring IoT complète et moderne, combinant :
\begin{itemize}
    \item \textbf{Surveillance Temps Réel :} Données instantanées du troupeau
    \item \textbf{Alertes Intelligentes :} Détection proactive des problèmes
    \item \textbf{Visualisations Riches :} Graphiques interactifs et métriques détaillées
    \item \textbf{Interface Moderne :} Design responsive et accessible
    \item \textbf{Performance Optimisée :} Expérience utilisateur fluide
\end{itemize}

Cette interface permet aux éleveurs de prendre des décisions éclairées basées sur des données fiables et actualisées en continu.

\end{document}
